% Generated by scripts/blueprint_gen.py; do not edit.
\chapter{The twisted theorem}\label{chap:Z3}
\textbf{Target.} For a compact Hausdorff complex manifold \texttt{X}, a \texttt{∂̄}-closed \texttt{h ∈ 𝒜\textasciicircum{}\{0,2\}(X)} and a twist datum \texttt{T} for \texttt{h}, the functor

\texttt{J\_T : Ho(P\_\{(𝒜\_X(X), ∂̄, h)\}) ⥤ D(TwMod O 𝒰 ȟ)}, \texttt{M ↦ ((α\textasciicircum{}h(M)|U\_α, 𝕏 − h⁰\_α), u\_αβ·)\_α}

is full and faithful with essential image the twisted perfect complexes \texttt{TwPerf}; and every closed \texttt{h} has a twist datum. Closest sources: [A, Thm 7.4.3] (first equivalence, compact \texttt{X}), [B, Thm 4.2.1] (stated without proof; \texttt{D\textasciicircum{}b\_coh} replaced by \texttt{D\_perf}, which is Z1).

\textbf{Shape of the proof.} On each \texttt{U\_α} the curved algebra is gauge-equivalent to the flat one (\texttt{𝕏 ↦ 𝕏 − h⁰\_α}), so every over-site lemma of C1/C2/C4/C7 and H2 on opens is reused on the cover pieces. Twisted sheaves have no global sections, so every global step (softness S3.b, globalization S10, lifting A11) acts on untwisted objects: Hom sheaves, a glued \texttt{𝒜⁰}-complex, global sections of a strictly perfect complex. New global work: the abelian category \texttt{TwMod} and its homological algebra, gluing of sheaves of modules (S12), twisted Hom sheaves, the glued curved model, A11 over the curved algebra (already general, §1.1).

\textbf{Size.} ≈ 12–21k lines (H2 was ≈ 6k); checked statement (statement closure, \texttt{exists\_twistData}, K8) ≈ 5–8k. Three stages, gates G1–G5, phases 8–10 (§5).

\begin{definition}[Z3.a — Curved dga/presheaf]
\label{Z3.a}
\lean{Cohesive.GCDga.withC, Cohesive.DgaPresheaf.withCurvature, Cohesive.DgaSheafData, Cohesive.PreResolvingData.toSheafData}
\leanok
\uses{C0}
Let A be a bounded graded-commutative curved dga (A1) with curvature c. For a closed element c′ ∈ A\textsuperscript{2} (dc′ = 0), A with curvature c′ is the curved dga with the same algebra, grading and differential and curvature c′. For a presheaf of dgas D on X and a closed h ∈ A(X)\textsuperscript{2}, the \emph{curved presheaf} D\textsubscript{h} has the same algebras, gradings and restriction maps, with curvature h|U on each open U. \emph{Sheaf data} on X are the fields of pre-resolving data (C0) except c = 0: a presheaf of possibly curved dgas with gluing and degree-locality of sections, 𝒜\textsuperscript{0} a fine sheaf of commutative rings, a structure sheaf O with ρ : O → 𝒜\textsuperscript{0} and d ∘ ρ = 0, and every 𝒜\textsuperscript{q} locally free over 𝒜\textsuperscript{0}. Every pre-resolving data are sheaf data, and sheaf data with D replaced by D\textsubscript{h} are again sheaf data.
\par\smallskip{\footnotesize Source: blueprint-tw.md §4.}
\end{definition}
\paragraph{Blueprint notes.}
\texttt{withC} (+ rfl API), \texttt{withCurvature}, \texttt{DgaSheafData}, \texttt{toSheafData}, bridges

Files: \texttt{Algebra/CurvedDga.lean}, \texttt{Comparison/CurvedPresheaf.lean}.

\begin{theorem}[Z3.b — Gauge]
\label{Z3.b}
\lean{Cohesive.CohesiveModule.gauge, Cohesive.QuasiCohesive.ungauge, Cohesive.gauge_hom, Cohesive.gauge_δ, Cohesive.gauge_Z0, Cohesive.gauge_B0}
\leanok
Let k be a commutative ring in which 2 is a unit, A a bounded graded-commutative curved dga over k with curvature c, β ∈ A\textsuperscript{1}, and c′ = c + dβ. For a (quasi-)cohesive module (Q, 𝕏) over (A, c′), its \emph{gauge} (Q, 𝕏 − β·) is a (quasi-)cohesive module over (A, c), with (𝕏 − β)\textsuperscript{2} = c′ − dβ + β\textsuperscript{2} = c; the inverse is 𝕏 ↦ 𝕏 + β·. The gauge does not change the Hom complexes: Hom\textsuperscript{n}(M₁, M₂) is the same submodule, δ agrees on it, hence so do the closed and exact degree-0 maps. Consequently the gauge is an equivalence Ho(P\textsubscript{(A, c + dβ)}) ≌ Ho(P\textsubscript{(A, c)}), the identity on morphisms.

Deviation: the sources write "2 ≠ 0"; Lean assumes 2 is a unit in k (true for k = ℂ).
\par\smallskip{\footnotesize Source: blueprint-tw.md §4.}
\end{theorem}
\begin{proof}
\leanok
\uses{Z3.a}
\textbf{A.2} gauge curvature.

\textbf{A.3} gauge on Hom: identity on \texttt{Hom\textasciicircum{}n}, \texttt{δ}, \texttt{Z0}, \texttt{B0}, composition.

\textbf{A.4} gauge commutes with base change.
\paragraph{Blueprint notes.}
A.2–A.4, \texttt{gaugeHoP} (explicit \texttt{@}-projections, P15)

Files: \texttt{Algebra/Gauge.lean}, \texttt{Algebra/GaugeBaseChange.lean}.
\end{proof}

\begin{theorem}[Z3.c — Curved Dolbeault, twist data]
\label{Z3.c}
\lean{Cohesive.DolbeaultClosed2, Cohesive.curvedDolbeaultPresheaf, Cohesive.curvedDolbeault, Cohesive.TwistCocycle, Cohesive.TwistCocycle.trivial, Cohesive.exactTwistDatum}
\leanok
Let X be a complex manifold modelled on a finite-dimensional complex normed space E.

\begin{itemize}
\item A \emph{closed (0,2)-form} is a global Dolbeault form h ∈ 𝒜\textsubscript{X}(X) of form degree 2 with ∂̄h = 0. The curved Dolbeault presheaf is U ↦ (𝒜\textsubscript{X}(U), ∂̄, h|U), and the curved Dolbeault algebra is (𝒜\textsubscript{X}(X), ∂̄, h), whose curvature is h.
\item A \emph{twist cocycle} for a sheaf of commutative rings R is a finite open cover (U\textsubscript{α}) of X with units ȟ\textsubscript{αβγ} ∈ R(U\textsubscript{α} ∩ U\textsubscript{β} ∩ U\textsubscript{γ})\textasciicircum{}×, equal to 1 when two indices coincide, satisfying ȟ\textsubscript{βγδ} ȟ\textsubscript{αγδ}\textsuperscript{−1} ȟ\textsubscript{αβδ} ȟ\textsubscript{αβγ}\textsuperscript{−1} = 1 on every open contained in U\textsubscript{α}, U\textsubscript{β}, U\textsubscript{γ}, U\textsubscript{δ}. The trivial cocycle is the cover \{X\} with ȟ = 1.
\item A \emph{twist datum} for a closed h ∈ 𝒜(X)\textsuperscript{2} over sheaf data S is a twist cocycle ȟ of S.O together with primitives h\textsuperscript{0}\textsubscript{α} ∈ 𝒜\textsuperscript{1}(U\textsubscript{α}), dh\textsuperscript{0}\textsubscript{α} = h|U\textsubscript{α}, and units u\textsubscript{αβ} ∈ 𝒜\textsuperscript{0}(U\textsubscript{α} ∩ U\textsubscript{β})\textasciicircum{}× with du\textsubscript{αβ} = (h\textsuperscript{0}\textsubscript{β} − h\textsuperscript{0}\textsubscript{α}) u\textsubscript{αβ}, u\textsubscript{αα} = 1, u\textsubscript{βα} = u\textsubscript{αβ}\textsuperscript{−1}, and ρ(ȟ\textsubscript{αβγ}) = u\textsubscript{βγ} u\textsubscript{αβ} u\textsubscript{αγ}\textsuperscript{−1}. For h = dβ with β ∈ 𝒜\textsuperscript{1}(X) there is a twist datum on the cover \{X\} with h\textsuperscript{0} = β, u = 1, ȟ = 1 (for β = 0: the trivial datum for h = 0).
\end{itemize}
Deviation: multiplicative data (h\textsuperscript{0}, u, ȟ) on a finite cover instead of Antweiler's (h\textsuperscript{0}, h\textsuperscript{1}, h\textsuperscript{2}, exp); the normalization and cocycle condition of ȟ are imposed as fields.
\par\smallskip{\footnotesize Source: blueprint-tw.md §4.}
\end{theorem}
\begin{proof}
\leanok
\uses{F14, Z3.a}
\begin{itemize}
\item \texttt{withC (i : GCDga 𝒜) (c') (hc : c' ∈ 𝒜 2) (hd : i.d c' = 0) : GCDga 𝒜 := \{ i with c := c', … \}}, \texttt{@[reducible]}, never an instance. \texttt{d}, \texttt{N}, \texttt{commRing0}, \texttt{Tot}, \texttt{Hom} agree by \texttt{rfl} (P15).
\item \texttt{DgaPresheaf.withCurvature D h hh hd}: same \texttt{A}, \texttt{𝒜}, grading; \texttt{dga U := withC (D.dga U) (D.res le\_top h) …}; \texttt{res} the same algebra maps with \texttt{map\_c} from \texttt{res\_comp} (built with explicit \texttt{@DgaHom.mk}, P15).
\item \texttt{DgaSheafData X} (u): \texttt{PreResolvingData} without \texttt{c\_eq\_zero}; \texttt{PreResolvingData.toSheafData}, \texttt{DgaSheafData.withCurvature}.
\item \textbf{Gauge} (A.2–A.4). For \texttt{M} over \texttt{withC i c'} and \texttt{β ∈ 𝒜¹} with \texttt{c' = c + dβ}: \texttt{gauge M β := (M.Q, 𝕏 − β·)} over \texttt{i}. \texttt{(𝕏 − β)² = c' − dβ + β² = c} by Leibniz \texttt{𝕏(βx) = dβ·x − β𝕏x}, and \texttt{β² = 0} since \texttt{β² = −β²} (graded commutativity). \warn{} Generic \texttt{k} needs no 2-torsion: hypothesis \texttt{IsUnit (2 : k)} in the generic gauge lemma (true for ℂ); deviation from A/B, which wrote "2 ≠ 0". Same \texttt{Hom\textasciicircum{}n}, same \texttt{δ} (the \texttt{β}-terms cancel since \texttt{φ(βx) = (−1)\textasciicircum{}n βφ(x)}), so \texttt{gaugeHoP : HoP (withC i c') ≌ HoP i}, identity on \texttt{H0}. \texttt{(gauge M β).baseChange f = gauge (M.baseChange f) (f β)} on generators. Dictionary: [B, §2.7] \texttt{(id, ω = β)}, sign reversed for left modules.
\item \texttt{DolbeaultClosed2 E X := \{h : DolbeaultAlg E ⊤ // h ∈ dolbeaultGrade E ⊤ 2 ∧ ∂̄h = 0\}}; \texttt{curvedDolbeaultPresheaf E X h := (dolbeaultPresheaf E X).withCurvature h.1 …}; \texttt{curvedDolbeault E X h := (curvedDolbeaultPresheaf E X h).𝒜 ⊤}.
\end{itemize}
\texttt{TwistCocycle R} (u): \texttt{ι : Type}, \texttt{[Fintype ι]}, \texttt{U : ι → Opens X}, \texttt{⨆ U = ⊤}, \texttt{ch α β γ : (R(U\_αβγ))ˣ}, normalized, Čech cocycle on quadruple intersections. \texttt{TwistDatum S h extends TwistCocycle S.O} (u): \texttt{h0 α ∈ 𝒜¹(U\_α)} with \texttt{d h0\_α = h|U\_α}; \texttt{u α β : 𝒜⁰(U\_αβ)ˣ} with \texttt{d u\_αβ = (h0\_β − h0\_α) u\_αβ}, \texttt{u\_αα = 1}, \texttt{u\_βα = u\_αβ⁻¹}; \texttt{ρ(ch\_αβγ) = u\_βγ u\_αβ u\_αγ⁻¹} on \texttt{U\_αβγ}. \texttt{TwistData E X h := TwistDatum (dolbeaultPreData E X) .toSheafData h.1}. Normalization and cocycle of \texttt{ch} follow from (4), (3) and injectivity of \texttt{ρ} (B correction 8). Special data: \texttt{trivialTwistDatum} (\texttt{h = 0}: \texttt{ι = Unit}, \texttt{h0 = 0}, \texttt{u = 1}, \texttt{ch = 1}), \texttt{exactTwistDatum h β} (\texttt{h = dβ}: \texttt{h0 = β}, \texttt{u = 1}, \texttt{ch = 1}).
\paragraph{Blueprint notes.}
§2.1–2.2, \texttt{trivialTwistDatum}, \texttt{exactTwistDatum}

Files: \texttt{Forms/CurvedDolbeault.lean}, \texttt{Sheaf/Twisted/CocycleDef.lean}, \texttt{Comparison/TwistDatumDef.lean}.
\end{proof}

\begin{theorem}[Z3.d — Exp, existence]
\label{Z3.d}
\lean{Cohesive.expForm, Cohesive.expUnit, Cohesive.dbar_expForm, Cohesive.exists_cover_primitives, Cohesive.twistData_of_primitives, Cohesive.exists_twistData}
\leanok
Let X be a compact Hausdorff complex manifold modelled on a finite-dimensional complex normed space E, and h ∈ 𝒜\textsuperscript{0,2}(X) with ∂̄h = 0. Then h has a twist datum. More precisely:

\begin{itemize}
\item there are a finite open cover (U\textsubscript{α}) of X, primitives h\textsuperscript{0}\textsubscript{α} ∈ 𝒜\textsuperscript{1}(U\textsubscript{α}) of h|U\textsubscript{α}, and functions η\textsubscript{αβ} ∈ 𝒜\textsuperscript{0}(U\textsubscript{α} ∩ U\textsubscript{β}) with ∂̄η\textsubscript{αβ} = h\textsuperscript{0}\textsubscript{β} − h\textsuperscript{0}\textsubscript{α}, η\textsubscript{αα} = 0 and η\textsubscript{βα} = −η\textsubscript{αβ};
\item for any such data, u\textsubscript{αβ} = exp(η\textsubscript{αβ}) (a unit of 𝒜\textsuperscript{0} with ∂̄ exp(η) = exp(η)·∂̄η) and the holomorphic units ȟ\textsubscript{αβγ} with ρ(ȟ\textsubscript{αβγ}) = u\textsubscript{βγ} u\textsubscript{αβ} u\textsubscript{αγ}\textsuperscript{−1} form a twist datum for h.
\end{itemize}
\par\smallskip{\footnotesize Source: blueprint-tw.md §4.}
\end{theorem}
\begin{proof}
\leanok
\uses{F12, H1, Z3.c, Z4.d, Z4.g, Z5.g}
\textbf{B.1} H1 (\texttt{q = 1}) + finite subcover: \texttt{U\_α}, \texttt{h0\_α}.

\textbf{B.2} \texttt{ω\_αβ := h0\_β − h0\_α} closed, antisymmetric.

\textbf{B.3} refinement: shrink (\texttt{exists\_iUnion\_eq\_closure\_subset}), at each \texttt{x} intersect H1 (\texttt{q = 0}) neighbourhoods in \texttt{U\_αβ} or complements of \texttt{cl U'\_α ∩ cl U'\_β}; finitely many \texttt{W\_i}; cover \texttt{U'\_α ⊓ W\_i} (index \texttt{ι × Fin m}): \texttt{ω} exact on all double intersections; symmetrized primitives \texttt{η'}.

\textbf{B.4} \texttt{u := exp ∘ η'} (\texttt{HasFDerivAt.cexp}, \texttt{ContDiff.cexp}; \texttt{∂̄} passes the holomorphic factor through, Z4.d \texttt{dbarPart\_comp\_clm}).

\textbf{B.5} \texttt{ζ = u\_bc u\_ab u\_ac⁻¹} is \texttt{∂̄}-closed, so \texttt{ζ = ρ(ch)} (F12, \texttt{ρ} injective).

\textbf{B.6} cocycle, normalization, cyclic symmetry \texttt{ch\_bca = ch\_abc} (used in C.4).
\paragraph{Blueprint notes.}
B.1–B.6

Files: \texttt{Forms/ExpForm.lean}, \texttt{Forms/TwistExists.lean}.
\end{proof}

\begin{theorem}[Z3.e — H2 on opens, LL\textsuperscript{h}]
\label{Z3.e}
\lean{Cohesive.curvedRes, Cohesive.curvedRestrict, Cohesive.curvedLL, Cohesive.localGaugeLemmaOpen, Cohesive.exists_free_baseChange_open}
\leanok
\uses{Z5.g, Z5.h, Z5.i}
Let X be a complex manifold modelled on a finite-dimensional complex normed space E.

\begin{itemize}
\item (H2 over an open.) For an open U₀ ⊆ X, a cohesive module N over 𝒜\textsubscript{X}(U₀) and x ∈ U₀, there are an open U′ ⊆ U₀ containing x, a bounded complex W of finitely generated projective O\textsubscript{X}(U′)-modules and an isomorphism N|U′ ≅ 𝒜\textsubscript{X}(U′) ⊗\textsubscript{O\_X(U′)} W in Ho(P\textsubscript{𝒜\_X(U′)}).
\item (Curved local lifting.) Let X moreover be compact Hausdorff, h ∈ 𝒜\textsuperscript{0,2}(X) closed, T a twist datum for h and M a cohesive module over (𝒜\textsubscript{X}(X), ∂̄, h). For x ∈ U\textsubscript{α} the same conclusion holds for the gauged restriction gauge\_\{h\textsuperscript{0}\textsubscript{α}\}(M|U\textsubscript{α}), a cohesive module over the flat algebra 𝒜\textsubscript{X}(U\textsubscript{α}) with operator 𝕏 − h\textsuperscript{0}\textsubscript{α}.
\end{itemize}
Deviation: Antweiler states the curved lemma for M|U′ itself; it is stated after the gauge by h\textsuperscript{0}\textsubscript{α}, which is an equivalence of homotopy categories (Z3.b).
\par\smallskip{\footnotesize Source: blueprint-tw.md §4.}
\end{theorem}
\begin{proof}
\leanok
\uses{Z3.b, Z3.c}
See the Lean proof and the cited items.
\paragraph{Blueprint notes.}
Z5.g–i with base \texttt{U₀}; \texttt{localGaugeLemma} re-derived at \texttt{⊤}; \texttt{curvedLL}

Files: \texttt{Forms/LocalFreeOpen.lean}, \texttt{Forms/LocalGaugeOpen.lean}, \texttt{Forms/CurvedLL.lean}.
\end{proof}

\begin{definition}[Z3.f — \texttt{restrOver}]
\label{Z3.f}
\lean{Cohesive.restrOver, Cohesive.restrOverAdj, Cohesive.restr_restrOver}
\leanok
\uses{S1}
Let R be a sheaf of rings on X and V ⊆ U opens. Restriction M ↦ M|V from sheaves of R|U-modules on the over-site of U to sheaves of R|V-modules on the over-site of V is Mathlib's pushforward along Over V → Over U. It is left adjoint to the pushforward along the other direction (which therefore preserves injective objects), it also has a left adjoint, so it is exact, and restricting to U and then to V is isomorphic to restricting to V.
\par\smallskip{\footnotesize Source: blueprint-tw.md §4.}
\end{definition}
\paragraph{Blueprint notes.}
P17 promoted: adjoints, exactness, \texttt{restr\_restrOver}, injectives

Files: \texttt{Sheaf/RestrOver.lean}.

\begin{definition}[Z3.g — \texttt{TwMod}]
\label{Z3.g}
\lean{Cohesive.sec, Cohesive.toSec, Cohesive.secRes, Cohesive.TwMod, Cohesive.TwMod.Hom, Cohesive.TwMod.comp}
\leanok
\uses{S1}
Let R be a sheaf of commutative rings on X, (U\textsubscript{α})\textsubscript{α∈ι} a family of opens and ȟ\textsubscript{αβγ} ∈ R(U\textsubscript{α} ∩ U\textsubscript{β} ∩ U\textsubscript{γ})\textasciicircum{}×. An \emph{ȟ-twisted sheaf of R-modules} consists of sheaves of modules M\textsubscript{α} on the over-sites of U\textsubscript{α}, and for every open W ⊆ U\textsubscript{α} ∩ U\textsubscript{β} an additive map φ\textsubscript{αβ} : M\textsubscript{α}(W) → M\textsubscript{β}(W), R(W)-linear and compatible with restriction, with φ\textsubscript{αα} = id and φ\textsubscript{βγ} ∘ φ\textsubscript{αβ} = ȟ\textsubscript{αβγ}|W · φ\textsubscript{αγ} on every W ⊆ U\textsubscript{α} ∩ U\textsubscript{β} ∩ U\textsubscript{γ}. Morphisms are families of morphisms M\textsubscript{α} → N\textsubscript{α} commuting with the gluing maps; composition is componentwise, and M ↦ M\textsubscript{α} is the component functor.

Deviation: the gluing is given on sections over every W ⊆ U\textsubscript{α} ∩ U\textsubscript{β} instead of as isomorphisms of restricted sheaves; φ\textsubscript{αα} = id is a field.
\par\smallskip{\footnotesize Source: blueprint-tw.md §4.}
\end{definition}
\paragraph{Blueprint notes.}
P16 promoted: \texttt{sec}, \texttt{toSec}, \texttt{secRes}, structure, \texttt{Hom}, \texttt{Category}, \texttt{comp}, \texttt{φ\_self}, \texttt{φ\_symm}

Files: \texttt{Sheaf/Twisted/Def.lean}.

\begin{theorem}[Z3.h — Abelian]
\label{Z3.h}
\lean{Cohesive.homOfSec, Cohesive.TwMod.Hom.glue_comm, Cohesive.TwMod.isIso_coimageImageComparison, Cohesive.TwMod.kerIsLimit, Cohesive.TwMod.cokIsColimit, Cohesive.sec_injective_of_mono}
\leanok
For every sheaf of commutative rings R on X, family of opens (U\textsubscript{α}) and family of units ȟ\textsubscript{αβγ} (no cocycle condition needed), the category of ȟ-twisted sheaves of R-modules is abelian. Kernels and cokernels are computed componentwise, with gluing maps induced from those of the source and target, and the component functors M ↦ M\textsubscript{α} preserve them. Supporting facts: a compatible family of R(W)-linear section maps over all W ⊆ V defines a morphism of restricted sheaves on V, and a monomorphism of sheaves of modules is injective on sections.
\par\smallskip{\footnotesize Source: blueprint-tw.md §4.}
\end{theorem}
\begin{proof}
\leanok
\uses{Z3.f, Z3.g}
\textbf{C.2} \texttt{Abelian}: componentwise kernels/cokernels (sections commute with finite limits; \texttt{restrOver} exact), \texttt{Abelian.ofCoimageImageComparisonIsIso} with \texttt{PreservesCoimageImageComparison.iso}; quasi-isomorphisms componentwise.

\textbf{C.2′} (refinement, 3 Nov 2026, P19): sections do not commute with cokernels (sheafification), so \texttt{φ} on a cokernel is not defined section by section. Bridge: a family \texttt{φ α β} over all \texttt{W ≤ U\_α ∩ U\_β}, \texttt{R(W)}-linear and natural, is the same as a morphism of sheaves \texttt{glueHom α β : restrOver inf\_le\_left (M α) ⟶ restrOver inf\_le\_right (M β)} on \texttt{Over (U\_α ∩ U\_β)} (sections of \texttt{restrOver} are the original sections by \texttt{rfl}, P19). Kernels: componentwise, \texttt{φ} restricted (sections commute with kernels). Cokernels: componentwise; \texttt{glueHom} is the map induced on cokernels (\texttt{restrOver} preserves cokernels); \texttt{φ\_self} and the cocycle are equalities of sheaf morphisms on \texttt{Over W}-type over-sites, checked after precomposing with the epimorphism \texttt{restrOver (coker.π)}; then read on sections. Implementation (3 Nov 2026): \texttt{homOfSec} turns a linear natural section family on \texttt{W' ≤ W} into a sheaf morphism on \texttt{Over W}; \texttt{sec\_eq\_of\_epi} cancels \texttt{restrOver p} (cocycle of a quotient, with the scalar \texttt{ȟ} as a section family, so no scalar endomorphism of sheaves is needed); sub-objects use injectivity of monos on sections; \texttt{TwMod.ofMono}/\texttt{ofEpi} package this once for kernels and cokernels; biproducts use joint injectivity of \texttt{fst, snd} on sections.

\textbf{C.3} products, coproducts componentwise.
\paragraph{Blueprint notes.}
C.2–C.3

Files: \texttt{Sheaf/Twisted/SecHom.lean}, \texttt{SubQuot.lean}, \texttt{Preadditive.lean}, \texttt{Kernels.lean}, \texttt{Cokernels.lean}, \texttt{Biprod.lean}, \texttt{Abelian.lean}.
\end{proof}

\begin{theorem}[Z3.i — Existence of injectives]
\label{Z3.i}
\lean{Cohesive.pushCounit, Cohesive.pushHomEquiv, Cohesive.TwistCocycle.unitRes_cyclic, Cohesive.TwMod.mono_of_mono_f, Cohesive.TwMod.comp_preservesMonomorphisms, Cohesive.push_preservesInjectiveObjects}
\leanok
Let R be a sheaf of commutative rings on X and ȟ a twist cocycle of R on a finite cover (U\textsubscript{α}).

\begin{itemize}
\item For each α, Toda's twisted pushforward N ↦ push\textsubscript{α}(N), from sheaves of R|U\textsubscript{α}-modules to ȟ-twisted sheaves (sections over W ⊆ U\textsubscript{β} are N(W ∩ U\textsubscript{α}), gluing by ȟ\textsubscript{αβγ}), is right adjoint to the component functor M ↦ M\textsubscript{α}.
\item The component functors preserve monomorphisms, and a morphism with mono components is mono; hence push\textsubscript{α} preserves injective objects.
\item Consequently the category of ȟ-twisted sheaves has enough injectives.
\item The cocycle satisfies the cyclic symmetry ȟ\textsubscript{βγα} = ȟ\textsubscript{αβγ}.
\end{itemize}
\par\smallskip{\footnotesize Source: blueprint-tw.md §4.}
\end{theorem}
\begin{proof}
\leanok
\uses{Z3.c, Z3.f, Z3.h}
\textbf{C.4} enough injectives: \texttt{push\_α N} (\texttt{β}-component \texttt{W ↦ N(W ⊓ U\_α)}, gluing \texttt{ch\_αβγ·}), \texttt{comp\_α ⊣ push\_α} (uses B.6), \texttt{push\_α} preserves injectives, \texttt{M ↪ ∏\_α push\_α I\_α} (components have enough injectives: \texttt{IsGrothendieckAbelian}).

\textbf{C.4′} (detail, 3 Nov 2026, P20). Fix a twist cocycle \texttt{T} (normalized, Čech cocycle). \textbf{(a)} Cyclic symmetry: the cocycle at \texttt{(α, β, γ, α)} with \texttt{ȟ\_αγα = ȟ\_αβα = 1} gives \texttt{ȟ\_βγα = ȟ\_αβγ} on every \texttt{W} in the three opens. \textbf{(b)} \texttt{j\_α N := overPullback (U\_α ≤ ⊤) N} on \texttt{X}; sections over \texttt{W} are \texttt{N(P\_α W)}, \texttt{P\_α W := pullback (W ≤ ⊤) (U\_α ≤ ⊤)} (a chosen pullback in \texttt{Opens X}, equal to \texttt{W ⊓ U\_α} but not definitionally; \texttt{P\_α W ≤ W}, \texttt{P\_α W ≤ U\_α}, \texttt{W ⊓ U\_α ≤ P\_α W}); restriction is restriction in \texttt{N} and \texttt{r ∈ R(W)} acts by \texttt{r|P\_α W} (P20, all \texttt{rfl}). \textbf{(c)} \texttt{push\_α N}: components \texttt{(push\_α N)\_β := restrOver (U\_β ≤ ⊤) (j\_α N)}, so sections over \texttt{W ≤ U\_β} are \texttt{N(P\_α W)} for every \texttt{β} (definitionally, independent of \texttt{β}: this is why one pushes to \texttt{X} first); gluing \texttt{φ\_βγ x := ȟ\_αβγ|P\_α W · x}. Linear (commutativity of \texttt{R}), natural (restriction of units), \texttt{φ\_ββ = id} (\texttt{ȟ\_αββ = 1}), cocycle \texttt{ȟ\_αγδ ȟ\_αβγ = ȟ\_βγδ ȟ\_αβδ} (Čech cocycle at \texttt{(α, β, γ, δ)}). On morphisms \texttt{push\_α k} has components \texttt{restrOver (j\_α k)} (commutes with scalars). \textbf{(d)} Counit \texttt{ε\_N : (push\_α N)\_α ⟶ N}, restriction along \texttt{W ≤ P\_α W} (\texttt{W ≤ U\_α}); an iso. \textbf{(e)} \texttt{Hom(M, push\_α N) ≃ Hom(M\_α, N)}: \texttt{F ↦ F\_α ≫ ε\_N}; inverse \texttt{g ↦ G}, \texttt{G\_β(x) := g(φ\_βα(x|P\_α W))} for \texttt{x ∈ M\_β(W)}. \texttt{G} commutes with the gluing: \texttt{G\_γ(φ\_βγ x) = g(φ\_γα φ\_βγ (x|P)) = g(ȟ\_βγα φ\_βα(x|P)) = ȟ\_αβγ G\_β(x)} by (a). Left inverse: \texttt{F.comm} at \texttt{(β, α)} and \texttt{ȟ\_αβα = 1} give \texttt{F\_α(φ\_βα y) = F\_β(y)}, so \texttt{G\_β x = ε(F\_β(x)|P)} and \texttt{P\_α(P\_α W) = P\_α W} up to restrictions (functoriality of restriction). Right inverse: \texttt{ε(G\_α x) = g(x)} by \texttt{φ\_αα = id} and naturality of \texttt{g}. Naturality in \texttt{M} is associativity, in \texttt{N} naturality of \texttt{ε}. Hence \texttt{comp\_α ⊣ push\_α} (\texttt{Adjunction.mkOfHomEquiv}). \textbf{(f)} \texttt{comp\_α} preserves monomorphisms (it preserves kernels), so \texttt{push\_α} preserves injectives (\texttt{Functor.preservesInjectiveObjects\_of\_adjunction\_of\_preservesMonomorphisms}). \textbf{(g)} Given \texttt{M}, choose monos \texttt{ι\_α : M\_α ↪ I\_α}, \texttt{I\_α} injective; \texttt{g := Pi.lift (fun α => (homEquiv)⁻¹ ι\_α) : M ⟶ ∏\_α push\_α I\_α} (finite product, injective). \texttt{g} is mono because each component \texttt{g.f β} is: \texttt{(g ≫ π\_β).f β ≫ ε = ι\_β} is mono. (A morphism of twisted sheaves with mono components is mono, by \texttt{hom\_ext}.) So \texttt{EnoughInjectives (TwMod R T.U T.ch)} for finite covers.

\textbf{C.5} restriction preserves injectives (exact left adjoints). (3 Nov 2026: for \texttt{comp\_α} the left adjoint is \texttt{extTw α}, built in Z3.l, so the \texttt{comp\_α} part of C.5 moves to Z3.l; \texttt{restrOver} part done in Z3.f.)
\paragraph{Blueprint notes.}
C.4–C.5

Files: \texttt{Sheaf/Twisted/Pushforward.lean}, \texttt{Injective.lean}.
\end{proof}

\begin{theorem}[Z3.j — Perfectness]
\label{Z3.j}
\lean{Cohesive.glueNatIso, Cohesive.IsPerfectOver, Cohesive.IsStrictlyPerfect.restrOver, Cohesive.TwIsPerfect, Cohesive.IsPerfectOverIn, Cohesive.IsPerfectOverIn.exactAt}
\leanok
Let R be a sheaf of commutative rings on X.

\begin{itemize}
\item A complex K of sheaves of R|V-modules on the over-site of an open V is \emph{perfect} if every x ∈ V has an open neighbourhood W ⊆ V with quasi-isomorphisms K|W ← L → P, P strictly perfect. If moreover all such P can be chosen concentrated in degrees [a, b], then the homology sheaves of K vanish outside [a, b]; restrictions of strictly perfect complexes are strictly perfect.
\item A complex K of ȟ-twisted sheaves is \emph{twisted perfect} if each component complex K\textsubscript{α} is perfect on the over-site of U\textsubscript{α}. D\textsubscript{perf}(X, ȟ) is the class of objects of the derived category of ȟ-twisted sheaves isomorphic to Q(K) for a twisted perfect K.
\item If ȟ\textsubscript{αβα} = 1 for all α, β, the gluing maps assemble into a natural isomorphism M\textsubscript{α}|U\textsubscript{α} ∩ U\textsubscript{β} ≅ M\textsubscript{β}|U\textsubscript{α} ∩ U\textsubscript{β}.
\end{itemize}
Deviation: twisted perfectness is componentwise (the gluing maps play no role), and local roofs of quasi-isomorphisms replace isomorphisms in the derived category.
\par\smallskip{\footnotesize Source: blueprint-tw.md §4.}
\end{theorem}
\begin{proof}
\leanok
\uses{S5, S9, Z3.c, Z3.h}
\textbf{C.6} \texttt{TwIsPerfect}, invariance, twisted S9.d (compactness). Detail of C.6 (3 Nov 2026, Z3.j). Restrictions compose definitionally: \texttt{restrOver h' ∘ restrOver h = restrOver (h'.trans h)} and \texttt{restrOver (le\_refl V) = 𝟭}, also on complexes (\texttt{Iso.refl} checks). \textbf{(a)} \texttt{restrOver} preserves summands of finite free modules (it commutes with \texttt{free}, as \texttt{restr} in S5.a) and strict bounds (additive), hence strict perfectness; it is exact, so preserves quasi-isomorphisms. \textbf{(b)} Over-site S9.c: \texttt{IsPerfectOver} is invariant under quasi-isomorphism (compose roofs; homotopy pullback of roofs through \texttt{D(R|W)}), verbatim as S9.c. \textbf{(c)} \texttt{IsPerfectOverIn K a b}: models in \texttt{[a, b]}. Restriction: a roof on \texttt{W ∋ x} restricts to \texttt{W ⊓ V'} by (a). Locality: if every \texttt{x ∈ V} has \texttt{V' ∋ x} with \texttt{IsPerfectOverIn (K|V') a b}, then \texttt{IsPerfectOverIn K a b} (models of \texttt{K|V'} are models of \texttt{K}, since \texttt{K|V'|W = K|W}). A single roof on \texttt{V} gives \texttt{IsPerfectOverIn K a b} (take \texttt{W = V}). \textbf{(d)} Homology of \texttt{K} with \texttt{IsPerfectOverIn K a b} vanishes outside \texttt{[a, b]}: verbatim S9.d, with locality of zero sheaves on \texttt{Over V} (sections over \texttt{W ≤ V} are determined on a cover of \texttt{W}). \textbf{(e)} \texttt{comp\_α} is exact (C.2), so it preserves quasi-isomorphisms and homology; a twisted sheaf with all components zero is zero. Hence \texttt{TwIsPerfect} is invariant under quasi-isomorphism and \texttt{TwPerf (Q K) ↔ TwIsPerfect K} (as S9.c), and \texttt{K.ExactAt i ↔ ∀ α, K\_α.ExactAt i}. \textbf{(f)} For a normalized \texttt{ch}, the gluing maps give an isomorphism of functors \texttt{comp\_α ⋙ resL α β ≅ comp\_β ⋙ resR α β} (\texttt{glueHom} is natural in \texttt{M} and has inverse built from \texttt{φ\_βα}, \texttt{φ\_symm}); so \texttt{K\_α|W ≅ K\_β|W} for every \texttt{W ≤ U\_α ⊓ U\_β}, as complexes. \textbf{(g)} Twisted S9.d: let \texttt{X} be compact, the cover cover \texttt{X}, \texttt{K} twisted perfect. For each \texttt{x} choose \texttt{β(x)} with \texttt{x ∈ U\_β(x)} and a roof for \texttt{K\_β(x)} on \texttt{W\_x ∋ x}, models in \texttt{[a\_x, b\_x]}; a finite subcover \texttt{W\_\{x\_1\}, …} gives one range \texttt{[a, b]}. For any \texttt{α} and \texttt{y ∈ U\_α}, \texttt{y ∈ W\_\{x\_i\}}; with \texttt{W' = W\_\{x\_i\} ⊓ U\_α}, the roof restricted to \texttt{W'} (c) and transported along \texttt{K\_β|W' ≅ K\_α|W'} (f) gives \texttt{IsPerfectOverIn (K\_α|W') a b}; locality (c) gives \texttt{IsPerfectOverIn K\_α a b}. So \texttt{TwIsPerfectIn K a b}, and by (d), (e) the homology of \texttt{K} vanishes outside \texttt{[a, b]}. \textbf{(h)} Sanity: \texttt{IsPerfectOver (K|⊤) ↔ IsPerfect K} (\texttt{restr\_restrOver}).
\paragraph{Blueprint notes.}
C.6

Files: \texttt{Sheaf/Twisted/PerfectDef.lean}, \texttt{PerfectOver.lean}, \texttt{Perfect.lean}.
\end{proof}

\begin{theorem}[Z3.k — Gluing]
\label{Z3.k}
\lean{Cohesive.glueEval, Cohesive.glueCoeval, Cohesive.ofSheafGlueIso, Cohesive.toGlueOfSheaf, Cohesive.rescale, Cohesive.isSheaf_gluePre}
\leanok
Let R be a sheaf of commutative rings on X and (U\textsubscript{α}) a family of opens.

\begin{itemize}
\item For a twisted sheaf M with trivial cocycle, the compatible families (s\textsubscript{α}) ∈ ∏\textsubscript{α} M\textsubscript{α}(W ∩ U\textsubscript{α}) form a sheaf glue(M) of R-modules on X, and M is isomorphic to the twisted sheaf of restrictions of glue(M) (no cover hypothesis).
\item If ⋃\textsubscript{α} U\textsubscript{α} = X, gluing is an equivalence between twisted sheaves with trivial cocycle and sheaves of R-modules on X, with inverse N ↦ (N|U\textsubscript{α}, identity gluing).
\item (Rescaling.) For units g\textsubscript{αβ} ∈ R(U\textsubscript{α} ∩ U\textsubscript{β})\textasciicircum{}× with g\textsubscript{αα} = 1 and families ȟ, ȟ′ with ȟ′\textsubscript{αβγ} g\textsubscript{αγ} = g\textsubscript{βγ} g\textsubscript{αβ} ȟ\textsubscript{αβγ}, replacing φ\textsubscript{αβ} by g\textsubscript{αβ} φ\textsubscript{αβ} is an equivalence between ȟ-twisted and ȟ′-twisted sheaves, the identity on components and morphisms.
\end{itemize}
\par\smallskip{\footnotesize Source: blueprint-tw.md §4.}
\end{theorem}
\begin{proof}
\leanok
\uses{S8, Z3.i}
\textbf{C.7} S12 gluing for \texttt{ch = 1}: \texttt{glue(N)(W) = \{(s\_α) : ψ\_αβ s\_α = s\_β\}}, \texttt{glue|U\_α ≅ N\_α} (cocycle only), \texttt{glueEquiv : TwMod R U 1 ≌ Mod R} (sheaf condition, \texttt{TopCat.Sheaf.existsUnique\_gluing}); \texttt{rescale g : TwMod R U ch ≌ TwMod R U (ch·ďg)}; tensor with summands of finite free \texttt{𝒜⁰}-modules commutes with \texttt{glue}. Detail of C.7 (3 Nov 2026, Z3.k; prototype P21). \textbf{(a)} Sections of a sheaf on \texttt{Over U} glue uniquely along open covers of \texttt{W ≤ U} (sheaf condition for the sieve generated by the cover, pullbacks are \texttt{⊓}). \textbf{(b)} For \texttt{M : TwMod R U 1}, \texttt{glue M (W) := \{(s\_α) ∈ ∏\_α M\_α(W ⊓ U\_α) : φ\_αβ(s\_α|W') = s\_β|W' for all W' ≤ W ⊓ U\_α ⊓ U\_β\}}, an \texttt{R(W)}-module via \texttt{R(W) → R(W ⊓ U\_α)}; restriction componentwise. (The condition is quantified over all \texttt{W'}, so no intersections are rewritten; by \texttt{φ\_res} it is the usual one.) Sheaf: given compatible families on a cover \texttt{W\_i} of \texttt{W}, for each \texttt{α} the components \texttt{s\textasciicircum{}i\_α} on \texttt{W\_i ⊓ U\_α} agree on overlaps, so glue (a) to \texttt{s\_α} on \texttt{W ⊓ U\_α}; the compatibility of \texttt{(s\_α)} holds locally on \texttt{W' ∩ W\_i}, hence by separatedness. \textbf{(c)} \texttt{ofSheaf (glue M) ≅ M}, no cover needed: on \texttt{W ≤ U\_γ}, \texttt{(s\_α) ↦ s\_γ|W} (\texttt{W ≤ W ⊓ U\_γ}), inverse \texttt{t ↦ (φ\_γα(t|W ⊓ U\_α))\_α}, compatible by the cocycle with \texttt{ȟ = 1} and \texttt{φ\_res}; inverse laws by \texttt{φ\_γγ = id} and compatibility of \texttt{(s\_α)} at \texttt{W' = W ⊓ U\_α}; it commutes with the gluing maps (cocycle again); natural in \texttt{M}. \textbf{(d)} For \texttt{⨆ U\_α = ⊤}, \texttt{glue (ofSheaf N) ≅ N}: \texttt{t ↦ (t|W ⊓ U\_α)\_α}, bijective by the sheaf condition of \texttt{N} for the cover \texttt{W ⊓ U\_α} of \texttt{W}; natural in \texttt{N}. \textbf{(e)} \texttt{glueEquiv : TwMod R U 1 ≌ Mod R} from (c), (d) (\texttt{Equivalence.mk}). \textbf{(f)} \texttt{rescale g}: for units \texttt{g\_αβ ∈ R(U\_αβ)ˣ} with \texttt{g\_αα = 1} and \texttt{ȟ'\_αβγ g\_αγ = g\_βγ g\_αβ ȟ\_αβγ} on every \texttt{W ≤ U\_αβγ}, \texttt{φ' := g\_αβ φ} defines \texttt{TwMod R U ȟ ⥤ TwMod R U ȟ'} (identity on components and morphisms); \texttt{g⁻¹} gives the inverse, unit and counit have identity components. With \texttt{g = 1} this is transport along \texttt{ȟ = ȟ'}. \textbf{(g)} "Tensor with summands of finite free \texttt{𝒜⁰}-modules commutes with \texttt{glue}" needs the twisted base change (Z3.l \texttt{BaseChange.lean}): moved to Z3.l.

\textbf{C.7′} (Z3.l′ detail, 5 Nov 2026, for review at G5; prototype P24). Fix a map of sheaves of commutative rings \texttt{ψ : R ⟶ S} on \texttt{X}; \texttt{S ⊗ −} is S8.a's \texttt{sheafBC} on an over-site, \texttt{η} the unit of \texttt{sheafBC ⊣ restrictScalars} (\texttt{η(m) = 1 ⊗ m}, as C7.b's \texttt{bcη}). \textbf{(a)} \emph{Base change commutes with restriction between over-sites}: for \texttt{V ≤ U} and \texttt{M} on \texttt{Over U}, \texttt{(S|U ⊗ M)|V ≅ S|V ⊗ M|V} (\texttt{restrOver}), natural in \texttt{M}, and on sections over \texttt{W ≤ V} it sends \texttt{η(m)} to \texttt{η(m)}. Proof: S8.a's \texttt{sheafBCOverIso} with \texttt{Over.map (V ≤ U)} in place of \texttt{Over.forget} (sectionwise the same presheaf; Mathlib's compatibility of pushforward along a continuous functor with sheafification, which applies to \texttt{Over.map} by P24). \textbf{(b)} \emph{Uniqueness}: two \texttt{S(W′)}-linear families on sections of \texttt{S|U ⊗ M} over \texttt{W′ ≤ W ≤ U}, natural in \texttt{W′}, with values in a sheaf of \texttt{S}-modules, are equal as soon as they agree on all \texttt{η(m)}, \texttt{m ∈ M(W′)}. Proof: as morphisms on \texttt{Over W} (\texttt{homOfSec}), transported by (a), they correspond under \texttt{sheafBC ⊣ restrictScalars} to their composites with \texttt{η}. \textbf{(c)} *\texttt{twBC ψ : TwMod R U ȟ ⥤ TwMod S U ψ(ȟ)}*: components \texttt{S|U\_α ⊗ M\_α}; gluing over \texttt{W ≤ U\_α ∩ U\_β} given by the sections of \texttt{(a)⁻¹ ∘ (S ⊗ g\_αβ) ∘ (a)}, where \texttt{g\_αβ} is \texttt{φ\_αβ} as a morphism on \texttt{Over (U\_α ∩ U\_β)} (\texttt{homOfSec}, C.2′). Linearity and naturality hold because this is a morphism of sheaves. Key formula: \texttt{φ′\_αβ(η m) = η(φ\_αβ m)} (from (a) and naturality of \texttt{η}). \texttt{φ′\_αα = id} and the cocycle \texttt{φ′\_βγ φ′\_αβ = ψ(ȟ\_αβγ)·φ′\_αγ} hold by (b): on \texttt{η(m)} both sides are \texttt{η(φ\_βγ φ\_αβ m) = η(ȟ\_αβγ·φ\_αγ m) = ψ(ȟ\_αβγ)·η(φ\_αγ m)} (\texttt{η} is \texttt{R}-linear into the restriction of scalars). On morphisms \texttt{S ⊗ f\_α} componentwise, which commutes with the gluing by (b). Functorial and additive (\texttt{sheafBC} is), and \texttt{comp\_α ∘ twBC = sheafBC ∘ comp\_α} by \texttt{rfl}. \textbf{(d)} \emph{Scope change}: the statement "tensor with summands of finite free \texttt{𝒜⁰}-modules commutes with \texttt{glue}" (C.7 (g)) follows from C7.a (\texttt{T} is a sheaf for every \texttt{F}) and C.7 (c). It is proved where it is used (Z3.t, G.5: assembling the global curved connection from local ones), not as a separate file of Z3.l′.
\paragraph{Blueprint notes.}
C.7 (a)–(f)

Files: \texttt{Sheaf/Twisted/GlueSec.lean}, \texttt{Glue.lean}, \texttt{GlueEquiv.lean}, \texttt{Rescale.lean}.
\end{proof}

\begin{theorem}[Z3.l — Generators]
\label{Z3.l}
\lean{Cohesive.extXRestrIso, Cohesive.extTwLiftSec, Cohesive.extTwLift, Cohesive.extTwUnit, Cohesive.extX_locally_zero, Cohesive.isFlat_extX}
\leanok
Let R be a sheaf of commutative rings on X and ȟ a twist cocycle of R on a finite cover (U\textsubscript{α}).

\begin{itemize}
\item For a sheaf F of R|V-modules, its extension by zero j\_!F is a sheaf of R-modules on X with (j\_!F)|V ≅ F, whose sections vanish near every point outside V; if F is flat, so is j\_!F.
\item The twisted extension by zero N ↦ ext\textsubscript{α}(N) (components (j\_!N)|U\textsubscript{β}, gluing by ȟ\textsubscript{αβγ}) is left adjoint to the component functor M ↦ M\textsubscript{α}.
\item Every twisted sheaf is a quotient of one with flat components, and every complex of twisted sheaves strictly zero above degree b receives a quasi-isomorphism from such a complex with componentwise flat terms.
\end{itemize}
\par\smallskip{\footnotesize Source: blueprint-tw.md §4.}
\end{theorem}
\begin{proof}
\leanok
\uses{S8, Z3.i, Z3.k}
Step C.5: see Z3.i.

\textbf{C.8} generators \texttt{extTw α (extZero\_V R)} with \texttt{Hom(G, M) ≅ M\_α(V)}; bounded-above resolutions by finite sums of them (\texttt{exists\_quasiIso\_of\_enough}), components flat. (v1.20, blueprint §11, approved 3 Nov 2026: explicit extension by zero \texttt{extX V F} on \texttt{X}, \texttt{extTw α N} with components \texttt{restr (U β) (extX (U α) N)} and gluing \texttt{ȟ\_αβγ ·}, \texttt{extTw α ⊣ comp\_α}, flatness by the S8.b argument, resolutions by \texttt{⊕\_α extTw α F\_α}; base change and C.7 (g) split off as Z3.l′.)
\paragraph{Blueprint notes.}
C.5 (\texttt{comp\_α}), C.8

Files: \texttt{Sheaf/Twisted/ExtZero.lean}, \texttt{ExtTw.lean}, \texttt{ExtTwAdj.lean}, \texttt{FlatResolution.lean}; \texttt{isFlat\_of\_locally\_zero} in \texttt{Sheaf/Flat.lean}.
\end{proof}

\begin{theorem}[Z3.l′ — Base change]
\label{Z3.lprime}
\lean{Cohesive.sheafBCRestrOverIso, Cohesive.sec_eq_of_bcUnit, Cohesive.bcφ_bcUnit}
\leanok
Let ψ : R → S be a morphism of sheaves of commutative rings on X.

\begin{itemize}
\item For V ⊆ U and a sheaf M of R|U-modules, base change commutes with restriction: (S|U ⊗ M)|V ≅ S|V ⊗ M|V.
\item Section maps out of S ⊗ M that are linear and natural are determined by their values on the elements 1 ⊗ m.
\item Componentwise base change is a functor from ȟ-twisted R-modules to ψ(ȟ)-twisted S-modules, whose gluing maps satisfy φ′\textsubscript{αβ}(1 ⊗ m) = 1 ⊗ φ\textsubscript{αβ}(m).
\end{itemize}
\par\smallskip{\footnotesize Source: blueprint-tw.md §4.}
\end{theorem}
\begin{proof}
\leanok
\uses{S8, Z3.h}
Step C.7: see Z3.k.

Step C.7′: see Z3.k.
\paragraph{Blueprint notes.}
C.7 (g), componentwise \texttt{sheafBC}; before Z3.u

Files: \texttt{Sheaf/Twisted/BaseChange.lean}, \texttt{GlueTensor.lean}.
\end{proof}

\begin{theorem}[Z3.m — C1 over \texttt{DgaSheafData}]
\label{Z3.m}
\lean{Cohesive.DgaSheafData.A0, Cohesive.DgaSheafData.𝒜X, Cohesive.DgaSheafData.r, Cohesive.DgaSheafData.R, Cohesive.DgaSheafData.TotSec, Cohesive.DgaSheafData.deg0}
\leanok
\uses{C1}
Let S be sheaf data on X (possibly curved), U₀ an open and M a cohesive module over 𝒜(U₀). The constructions of C1 are defined over S word for word: the sheaf of rings 𝒜\textsuperscript{0}, the global algebra 𝒜(X) (with the curvature of S), the restrictions r\textsubscript{V}, the ring sheaf O, the total module Tot(M|V) for V ⊆ U₀, and its elements of total degree n, an O(V)-module α(M)\textsuperscript{n}(V). These form a sheaf α(M)\textsuperscript{n} of O-modules with the operator 𝕏 : α(M)\textsuperscript{n} → α(M)\textsuperscript{n+1}, and α acts on closed degree-0 maps and homotopies. For pre-resolving data these agree definitionally with C1.

Deviation: no complex α(M) and no functor are built here, since 𝕏\textsuperscript{2} = c· need not vanish.
\par\smallskip{\footnotesize Source: blueprint-tw.md §4.}
\end{theorem}
\begin{proof}
\leanok
\uses{Z3.a}
\textbf{D.6} local identification \texttt{(J\_T M)\_α|W' ≅ sheafOfOver\_\{W'\}(gauge\_\{h0\_α\}(M|U\_α)|W')} (C1.c, A.2/A.4).
\paragraph{Blueprint notes.}
copies + rfl bridges (done 3 Nov 2026: AlphaDef, C1.a, C1.b section and degree maps; no complex/\texttt{J}, they need \texttt{c = 0}; \texttt{AlphaOver} not copied, D.6 applies it to the flat gauged module)

Files: \texttt{Comparison/SheafDataAlphaDef.lean}, \texttt{SheafDataAlpha.lean}, \texttt{SheafDataMap.lean}.
\end{proof}

\begin{theorem}[Z3.n — \texttt{J\_T}]
\label{Z3.n}
\lean{Cohesive.DgaSheafData.smul0L, Cohesive.DgaSheafData.gaugeDV, Cohesive.DgaSheafData.gaugeDV_gaugeDV, Cohesive.DgaSheafData.gaugeDV_smul0L, Cohesive.DgaSheafData.bcHomL_gauge_homotopy, Cohesive.TwistDatum.twTerm}
\leanok
Let S be sheaf data on X, h ∈ 𝒜(X)\textsuperscript{2} closed, T a twist datum for h over S, and let M be a cohesive module over the curved algebra (𝒜(X), h).

\begin{itemize}
\item For any sheaf data, β ∈ 𝒜\textsuperscript{1}(V) and M over 𝒜(U₀), V ⊆ U₀, the gauged operator 𝕏 − β· on α(M)(V) satisfies (𝕏 − β)\textsuperscript{2} = 0 if dβ is the curvature over V (h|V for the curved data); if du = (β′ − β)u for u ∈ 𝒜\textsuperscript{0}(V), then (𝕏 − β′)(u·x) = u·(𝕏 − β)x; and homotopies of maps give homotopies for the gauged operators.
\item J\textsubscript{T}(M) is the complex of ȟ-twisted O-modules with components J\textsubscript{T}(M)\textsuperscript{n\_α} = α(M)\textsuperscript{n}|U\textsubscript{α}, gluing maps u\textsubscript{αβ}· and differential 𝕏 − h\textsuperscript{0}\textsubscript{α} on U\textsubscript{α}.
\item M ↦ J\textsubscript{T}(M) defines functors J\textsubscript{T} : Ho(P\textsubscript{(𝒜(X), h)}) ⥤ K(TwMod(O, U, ȟ)) and, for any derived category structure, Ho(P\textsubscript{(𝒜(X), h)}) ⥤ D(TwMod(O, U, ȟ)).
\end{itemize}
Deviation: the gluing is by the units u\textsubscript{αβ} instead of Antweiler's exp(h\textsuperscript{1}\textsubscript{αβ}).
\par\smallskip{\footnotesize Source: blueprint-tw.md §4.}
\end{theorem}
\begin{proof}
\leanok
\uses{C1, Z3.i, Z3.m}
\textbf{D.1} terms (curved C1.a copy).

\textbf{D.2} \texttt{d\_α = 𝕏 − h0\_α}, \texttt{d\_α² = h − dh0\_α = 0}.

\textbf{D.3} gluing \texttt{u\_αβ·}, chain map by A.5, cocycle by (4).

\textbf{D.4} morphisms and homotopies (C1, \texttt{k} commutes with \texttt{u}).
\paragraph{Blueprint notes.}
D.1–D.4

Files: \texttt{Comparison/TwistedFunctorDef.lean}, \texttt{TwistedFunctor.lean}.
\end{proof}

\begin{theorem}[Z3.o — Local identification]
\label{Z3.o}
\lean{Cohesive.twLocalIso, Cohesive.curvedResM, Cohesive.gaugedRes, Cohesive.twCmp, Cohesive.twCmp_X}
\leanok
Let P be pre-resolving data on X, h ∈ 𝒜(X)\textsuperscript{2} closed, T a twist datum for h over P and M a cohesive module over (𝒜(X), h). For each α, the restriction M|U\textsubscript{α} is a cohesive module over (𝒜(U\textsubscript{α}), h|U\textsubscript{α}), and its gauge gauge\_\{h\textsuperscript{0}\textsubscript{α}\}(M|U\textsubscript{α}) (Z3.b, h|U\textsubscript{α} = dh\textsuperscript{0}\textsubscript{α}) is a cohesive module over the flat algebra 𝒜(U\textsubscript{α}). The comparison Tot(M|V) → Tot(gauge\_\{h\textsuperscript{0}\textsubscript{α}\}(M|U\textsubscript{α})|V), V ⊆ U\textsubscript{α}, intertwines 𝕏 − h\textsuperscript{0}\textsubscript{α} with the operator of the gauged module, and the component complex J\textsubscript{T}(M)\textsubscript{α} on the over-site of U\textsubscript{α} is isomorphic to α\textsubscript{U\_α}(gauge\_\{h\textsuperscript{0}\textsubscript{α}\}(M|U\textsubscript{α})) (C1).
\par\smallskip{\footnotesize Source: blueprint-tw.md §4.}
\end{theorem}
\begin{proof}
\leanok
\uses{Z3.b, Z3.m, Z3.n}
Step D.6: see Z3.m.
\paragraph{Blueprint notes.}
D.6 (generic \texttt{P : PreResolvingData}; \texttt{gaugedRes}, \texttt{twCmp})

Files: \texttt{Comparison/TwistedLocal.lean}; \texttt{Scratch/P23\_TwLocal.lean}.
\end{proof}

\begin{theorem}[Z3.p — Image of J]
\label{Z3.p}
\lean{Cohesive.PreResolvingData.restrOverSheafOfIso, Cohesive.isPerfectOver_sheafOfOver, Cohesive.twIsPerfect_twistedSheafOf}
\leanok
\uses{C1, H2}
Let X be a compact Hausdorff complex manifold modelled on a finite-dimensional complex normed space E.

\begin{itemize}
\item For pre-resolving data, opens U′ ⊆ U and a cohesive module N over 𝒜(U), α\textsubscript{U}(N)|U′ ≅ α\textsubscript{U′}(N|U′) as complexes on the over-site of U′.
\item For every open U and every cohesive module N over 𝒜\textsubscript{X}(U), α\textsubscript{U}(N) is perfect on U.
\item For a closed h ∈ 𝒜\textsuperscript{0,2}(X), a twist datum T for h and a cohesive module M over (𝒜\textsubscript{X}(X), ∂̄, h), the complex J\textsubscript{T}(M) is twisted perfect.
\end{itemize}
\par\smallskip{\footnotesize Source: blueprint-tw.md §4.}
\end{theorem}
\begin{proof}
\leanok
\uses{C2, Z3.e, Z3.j, Z3.o, Z4.g}
\textbf{E.1} LL\textsuperscript{h} = H2 on opens (Z3.e) for \texttt{gauge\_\{h0\_α\}(M|U\_α)}.

\textbf{E.2} C2's over-site core (base generalized \texttt{⊤ → U₀}, V3) + D.6.
\paragraph{Blueprint notes.}
E.1–E.2 (H2 over opens for any \texttt{N} over \texttt{𝒜(U)}, so no \texttt{curvedLL} bridge is needed)

Files: \texttt{Comparison/AlphaOverRestr.lean} (C1.c between opens), \texttt{Comparison/TwistedPerfect.lean}.
\end{proof}

\begin{theorem}[Z3.q — C3 curved]
\label{Z3.q}
\lean{Cohesive.DgaSheafData.ℋ}
\leanok
Let S be sheaf data on X (possibly curved, curvature c) and M₁, M₂ cohesive modules over 𝒜(X). The \emph{Hom complex of sheaves} ℋ(M₁, M₂) is the complex of sheaves of O-modules whose degree-n sections over V are the degree-n maps χ : Tot(M₁) → Tot(M₂|V) with χ(a·x) = (−1)\textsuperscript{n|a|} a|V·χ(x), with differential δχ = 𝕏₂ ∘ χ − (−1)\textsuperscript{n} χ ∘ 𝕏₁; δ\textsuperscript{2} = 0 because 𝕏₁\textsuperscript{2} and 𝕏₂\textsuperscript{2} are both multiplication by c. Its global sections are Hom\textsuperscript{•\_}\{P\textsubscript{A}\}(M₁, M₂) (a bijection compatible with δ), and if X is compact Hausdorff its terms are soft.
\par\smallskip{\footnotesize Source: blueprint-tw.md §4.}
\end{theorem}
\begin{proof}
\leanok
\textbf{F.2} curved ℋ\textsuperscript{h} (C3 over \texttt{DgaSheafData}); \texttt{𝓗om\textasciicircum{}tw(K, I)} = glued over-site Hom sheaves (conjugation, \texttt{ch} cancels); sections = \texttt{HomComplex} cochains of \texttt{TwMod}; \texttt{Ψ\textasciicircum{}tw = r\_α ∘ sheafOfMapV}.
\paragraph{Blueprint notes.}
F.2 (ℋ\textsuperscript{h})

Files: \texttt{Comparison/CurvedHomSheaf*.lean}, \texttt{CurvedHomGauge.lean}.
\end{proof}

\begin{theorem}[Z3.r — \texttt{𝓗om\textasciicircum{}tw}]
\label{Z3.r}
\lean{Cohesive.secHomOf, Cohesive.twHomComplex, Cohesive.twHomD_restr, Cohesive.hom_ext_of_restr, Cohesive.twHomRestrIso, Cohesive.twHomRestrOverIso}
\leanok
Let R be a sheaf of commutative rings on X, (U\textsubscript{α}) a cover of X, ȟ a family of units with ȟ\textsubscript{αβα} = 1, and K, L complexes of ȟ-twisted sheaves.

\begin{itemize}
\item For twisted sheaves K, L, the Hom sheaves Hom(K\textsubscript{α}, L\textsubscript{α}) with gluing s ↦ φ\textsuperscript{L\_αβ} ∘ s ∘ φ\textsuperscript{K\_βα} form a twisted sheaf with trivial cocycle; its gluing has global sections Hom(K, L).
\item 𝓗om\textsuperscript{tw}(K, L) is the complex of sheaves of R-modules on X with 𝓗om\textsuperscript{tw}(K, L)\textsuperscript{n} = ∏\textsubscript{p+n=q} glue(Hom(K\textsuperscript{p}, L\textsuperscript{q})), Mathlib's sign for the differential.
\item For V ⊆ U\textsubscript{α}, 𝓗om\textsuperscript{tw}(K, L)|V ≅ 𝓗om(K\textsubscript{α}|V, L\textsubscript{α}|V) as complexes on the over-site of V, and morphisms of sheaves on X agreeing on every U\textsubscript{α} agree.
\end{itemize}
\par\smallskip{\footnotesize Source: blueprint-tw.md §4.}
\end{theorem}
\begin{proof}
\leanok
\uses{S7, Z3.k}
Step F.2: see Z3.q.
\paragraph{Blueprint notes.}
F.2

Files: \texttt{Sheaf/Twisted/HomSheaf.lean}, \texttt{HomSections.lean}.
\end{proof}

\begin{theorem}[Z3.s — Full faithfulness]
\label{Z3.s}
\lean{Cohesive.bijective_twistedJ_map, Cohesive.full_twistedJ_dolbeault, Cohesive.faithful_twistedJ_dolbeault, Cohesive.TwistDatum.Ψtw, Cohesive.TwistDatum.Ψtw_top, Cohesive.LLOpen, Cohesive.jKtw, Cohesive.evLtw, Cohesive.ιtw}
\leanok
Let D be resolving data (C0) on a compact Hausdorff space X satisfying the local gauge lemma over every open (LL\textsubscript{open}: each cohesive module over 𝒜(U₀) is near each x ∈ U₀ isomorphic to 𝒜(U′) ⊗\textsubscript{O(U′)} W, W bounded finitely generated projective). Let h ∈ 𝒜(X)\textsuperscript{2} be closed and T a twist datum for h. Then:

\begin{itemize}
\item there is a morphism of complexes Ψ\textsuperscript{tw} : ℋ\textsuperscript{h}(M₁, M₂) → 𝓗om\textsuperscript{tw}(J\textsubscript{T} M₁, J\textsubscript{T} M₂), which on global sections sends a closed degree-0 map φ to J\textsubscript{T}(φ), and becomes a quasi-isomorphism after composing with a bounded below injective resolution of J\textsubscript{T}(M₂);
\item J\textsubscript{T} : Ho(P\textsubscript{(𝒜(X), h)}) ⥤ D(TwMod(O, U, ȟ)) is bijective on morphisms.
\end{itemize}
For the Dolbeault data of a compact Hausdorff complex manifold LL\textsubscript{open} holds, so J\textsubscript{T} is full and faithful for every closed h ∈ 𝒜\textsuperscript{0,2}(X) and every twist datum T.
\par\smallskip{\footnotesize Source: blueprint-tw.md §4.}
\end{theorem}
\begin{proof}
\leanok
\uses{A10, C2, C3, C4, S3.b, S4, S5, S6, S7, Z3.e, Z3.p, Z3.r}
\textbf{F.1} injective resolution \texttt{r : J\_T M₂ → I} (C.4).

Step F.2: see Z3.q.

\textbf{F.3} locally \texttt{Ψ\textasciicircum{}tw} is C4's \texttt{Ψ\_\{U'\}} for the gauged flat modules (A.3, C.7, D.6), quasi-iso by C4 step 3 + LL\textsuperscript{h}; globalize (S6).

\textbf{F.4} softness (S7.a′, S3.b).

\textbf{F.5} \texttt{H⁰} dictionary (\texttt{Qh\_map\_bijective\_of\_isKInjective}, \texttt{isKInjective\_of\_injective}). \textbf{F′} (detail, 4 Nov 2026, approved at G4 5 Nov 2026; refines F.1–F.5, no decision changed). Fix sheaf data, \texttt{h}, a twist datum \texttt{T} (normalized \texttt{ȟ}), \texttt{S\_h} the curved data, \texttt{M₁, M₂} cohesive over \texttt{(𝒜(X), h)}. \textbf{(a)} \emph{Curved C3} (Z3.q): C3's section data over \texttt{S : DgaSheafData} word for word (\texttt{HomSec}, \texttt{homRes}, \texttt{homδV}, \texttt{ℋ}, \texttt{homTop}). Only \texttt{δ² = 0} changes: \texttt{δ²χ = 𝕏₂²χ − χ𝕏₁² = c·χ − χ(c·) = 0} (\texttt{c} of degree 2, so \texttt{χ(c·y) = c·χ(y)}). \texttt{ℋ\textasciicircum{}h(M₁, M₂)(X) ≅ Hom\textasciicircum{}•(M₁, M₂)} as in C3. Softness as C3.a(ii). \textbf{(b)} \emph{Hom sections on over-sites.} For \texttt{P, Q} on \texttt{Over U} and \texttt{W ≤ U}, morphisms \texttt{P.over (W → U) ⟶ Q.over (W → U)} (sections of \texttt{homSheaf P Q} over \texttt{W}) are the same as \texttt{R(W′)}-linear natural section families on \texttt{W′ ≤ W} (\texttt{homOfSec₀}-type bridge, P19 rfl sections). \textbf{(c)} *\texttt{homTw K L}* (Z3.r), for \texttt{K, L : TwMod R U ȟ}: the untwisted (\texttt{ȟ = 1}) twisted sheaf with components \texttt{homSheaf (K\_α) (L\_α)} on \texttt{Over U\_α} and gluing by conjugation on sections over \texttt{W ≤ U\_α ∩ U\_β}: \texttt{s ↦ φ\textasciicircum{}L\_αβ ∘ s ∘ φ\textasciicircum{}K\_βα} (by (b), on \texttt{W′ ≤ W}). Linear (target action commutes with \texttt{φ\textasciicircum{}L}), natural, \texttt{φ\_self} from \texttt{φ\textasciicircum{}K\_αα = φ\textasciicircum{}L\_αα = id}; cocycle: \texttt{φ\textasciicircum{}K\_βα φ\textasciicircum{}K\_γβ = ȟ\_αβγ⁻¹ φ\textasciicircum{}K\_γα} (inverse of the cocycle, \texttt{φ\_symm}) and \texttt{φ\textasciicircum{}L\_βγ φ\textasciicircum{}L\_αβ = ȟ\_αβγ φ\textasciicircum{}L\_αγ}, so the scalars cancel (\texttt{s} is \texttt{R}-linear). Functorial: \texttt{homTwPost g}, \texttt{homTwPre f} (componentwise \texttt{homPost}, \texttt{homPre}). \textbf{(d)} *\texttt{𝓗om\textasciicircum{}tw(K, L)}* for complexes of \texttt{TwMod R U ȟ}: degree \texttt{n} is \texttt{∏\_\{p + n = q\} glue(homTw(Kᵖ, L\textasciicircum{}q))} in \texttt{Mod R} (\texttt{glue}, C.7 (b); \texttt{⨆ U = ⊤}), differential with Mathlib's sign as S7 (\texttt{homComplexD}), \texttt{δ² = 0} formally from functoriality of \texttt{homTwPost}/\texttt{homTwPre}; functorial in both variables (\texttt{twHomMap₂}, \texttt{twHomMap₁}). \textbf{(e)} \emph{Global sections}: \texttt{Γ(𝓗om\textasciicircum{}tw(K, L))ⁿ ≅ Cochain\_\{TwMod\}(K, L, n)} compatibly with \texttt{δ} and with \texttt{twHomMap₂} (post-composition): a global section of \texttt{glue(homTw(Kᵖ, L\textasciicircum{}q))} is a family of morphisms \texttt{Kᵖ\_α → L\textasciicircum{}q\_α} compatible with conjugation, i.e. a morphism of twisted sheaves. \textbf{(f)} \emph{Restriction}: for \texttt{V ≤ U\_α}, \texttt{𝓗om\textasciicircum{}tw(K, L)|V ≅ 𝓗om\_\{Over V\}(K\_α|V, L\_α|V)} as complexes (\texttt{glue(N)|U\_α ≅ N\_α}, C.7 (c); restriction preserves products; \texttt{restrOver} commutes with \texttt{homSheaf}), natural in \texttt{L} (and in \texttt{K}). \textbf{(g)} \emph{Twisted S5.c}: for \texttt{I} injective in \texttt{TwMod}, every restriction \texttt{(I\_α)|V} is injective (\texttt{I} is a summand of \texttt{∏\_β push\_β J\_β}, C.4 (g); the components of \texttt{push\_β J} are restrictions of the injective \texttt{j\_β J} on \texttt{X}, \texttt{overPullback} preserves injectives; \texttt{restr\_restrOver}; S5.c). \textbf{(h)} \emph{Softness, bounds}: \texttt{𝒜⁰(X)} acts on each \texttt{J\_T(M₁)\textasciicircum{}p} by morphisms of twisted sheaves (multiplication commutes with \texttt{u\_αβ·}), hence on the terms of \texttt{𝓗om\textasciicircum{}tw(J\_T M₁, I)} by precomposition, unitally and vanishing over \texttt{V} where \texttt{a|V = 0} (S3.b(i)); bounded below for \texttt{I} bounded below (\texttt{J\_T M₁} is bounded above). \textbf{(i)} *\texttt{Ψ\textasciicircum{}tw : ℋ\textasciicircum{}h(M₁, M₂) ⟶ 𝓗om\textasciicircum{}tw(J\_T M₁, J\_T M₂)}*: \texttt{χ ∈ ℋ\textasciicircum{}\{h,n\}(V)} goes, in component \texttt{α}, over \texttt{W ≤ V ⊓ U\_α}, degree \texttt{p}, to C4's \texttt{extAlong (χ|W)}. It commutes with the gluing (\texttt{u\_αβ} has degree \texttt{0}) and with \texttt{δ} (\texttt{h⁰\_α} has degree 1: \texttt{χ(h⁰\_α y) = (−1)ⁿ h⁰\_α χ(y)}, so the \texttt{h⁰\_α}-terms cancel). \textbf{(j)} \emph{Local} (F.3): \texttt{x ∈ U\_α}; LL over opens (Z3.e) for \texttt{gM₁ := gauge\_\{h⁰\_α\}(M₁|U\_α)} gives \texttt{U′ ∋ x}, \texttt{U′ ≤ U\_α}, \texttt{W} and a homotopy equivalence \texttt{φ : 𝒜(U′) ⊗ W → gM₁|U′}; with C2's \texttt{L, ε, u} (C4 step 3), \texttt{ι := u ≫ α\_\{U′\}(φ) ≫ (α\_\{U\_α\}(gM₁)|U′ ≅ α\_\{U′\}(gM₁|U′))⁻¹ ≫ (twLocalIso|U′)⁻¹ : L → (J\_T M₁)\_α|U′}. \texttt{Φ\textasciicircum{}tw := Ψ\textasciicircum{}tw|U′ ≫ (f) ≫ ι\textasciicircum{}*} is a sectionwise quasi-isomorphism: over \texttt{V ≤ U′} it is \texttt{ℋ\textasciicircum{}h(V) ≅ Hom\_along(gM₁|U′, gM₂|V)} (gauge is the identity on \texttt{Hom}, A.3; base change, as C4's \texttt{jK}) \texttt{→φ\textasciicircum{}*→ Hom\_along(𝒜(U′) ⊗ W, gM₂|V) ≅ Hom\_\{O(U′)\}(W, Tot(gM₂|V))} (A10.a) \texttt{≅ Cochain(L|V, α\_\{U\_α\}(gM₂)|V)} (S4.c, C4's \texttt{evL} with base \texttt{U\_α}) \texttt{≅ Cochain(L|V, (J\_T M₂)\_α|V)} (D.6). Then \texttt{r\_*} (S7.b′; \texttt{r\_α|U′} a quasi-iso, \texttt{comp\_α}, \texttt{restrOver} exact) and \texttt{ι\textasciicircum{}*} (S7.c′, (g)): \texttt{(Ψ\textasciicircum{}tw ≫ r\_*)|U′} is a quasi-isomorphism; S6. \textbf{(k)} \emph{Global} (F.4–F.5): \texttt{Γ(Ψ\textasciicircum{}tw ≫ r\_*)} is a quasi-isomorphism (S3.b(iv), (h)); \texttt{H⁰Γ ℋ\textasciicircum{}h = Hom\_\{Ho(P)\}(M₁, M₂)} ((a)), \texttt{H⁰Γ 𝓗om\textasciicircum{}tw(J\_T M₁, I) = K(TwMod)(J\_T M₁, I) = D(TwMod)(J\_T M₁, I)} ((e); \texttt{I} K-injective); the composite is \texttt{f ↦ J\_T(f) ≫ Q(r)}. Stated for resolving data \texttt{D} (compact Hausdorff \texttt{X}) with LL over opens as a hypothesis (\texttt{LLOpen}); for the Dolbeault data it is proved (\texttt{llOpen\_dolbeault}, from Z3.e), so the Dolbeault theorems carry no such hypothesis.
\paragraph{Blueprint notes.}
F.1–F.5

Files: \texttt{Comparison/TwistedFullyFaithful*.lean}.
\end{proof}

\begin{theorem}[Z3.t — Curved model]
\label{Z3.t}
\lean{Cohesive.twModel_TD_comm, Cohesive.isRescale_untwist, Cohesive.untwistRestrIso, Cohesive.untwist_compat, Cohesive.untwist_ext, Cohesive.πTot}
\leanok
Let P be pre-resolving data on X, h ∈ 𝒜(X)\textsuperscript{2} closed, T a twist datum for h and K′ a complex of ȟ-twisted O-modules.

\begin{itemize}
\item The cocycle ρ(ȟ) of 𝒜\textsuperscript{0}-units, rescaled by u\textsuperscript{−1}, is trivial; hence untwist(M) = glue(rescale\_\{u\textsuperscript{−1}\}(𝒜\textsuperscript{0} ⊗ M)) is a sheaf of 𝒜\textsuperscript{0}-modules on X with untwist(M)|U\textsubscript{α} ≅ 𝒜\textsuperscript{0}|U\textsubscript{α} ⊗ M\textsubscript{α}, whose sections are determined by their components, which satisfy s\textsubscript{β} = u\textsubscript{αβ}\textsuperscript{−1}·φ\textsubscript{αβ}(s\textsubscript{α}).
\item For X compact and connections θ\textsubscript{α} on 𝒜\textsuperscript{0}|U\textsubscript{α} ⊗ K′\textsubscript{α} (C7), the differentials of the local models commute with the gluing, giving a twisted model complex of K′.
\item For F = untwist(K′), the maps Tot(𝒜(X), F(X)) → Tot(𝒜(U\textsubscript{α}), F\textsubscript{α}(U\textsubscript{α})) jointly determine elements, and there is a curved operator θ\textsuperscript{h} with (θ\textsuperscript{h})\textsuperscript{2} = h restricting to θ\textsubscript{α} + h\textsuperscript{0}\textsubscript{α}.
\end{itemize}
\par\smallskip{\footnotesize Source: blueprint-tw.md §4.}
\end{theorem}
\begin{proof}
\leanok
\uses{C7, Z3.k, Z3.lprime, Z3.o}
\textbf{G.5} curved model: \texttt{∇\textasciicircum{}h\_α := ∇\_α + h0\_α ⊗ −} glue to \texttt{∇\textasciicircum{}h} with \texttt{(∇\textasciicircum{}h)² = h}; \texttt{N := (Γ(F\textasciicircum{}•), 𝕏\textasciicircum{}h)} quasi-cohesive over \texttt{(𝒜(X), h)} (curved C7.c copy).

\textbf{G.7} \texttt{e\_α ∈ Z0(gauge(M|U\_α), N\_α)}.
\paragraph{Blueprint notes.}
G.5, G.7

Files: \texttt{Comparison/CurvedConn.lean}, \texttt{TwistedModel.lean}.
\end{proof}

\begin{theorem}[Z3.u — Essential surjectivity]
\label{Z3.u}
\lean{Cohesive.exists_twistedJ_obj_iso, Cohesive.TwIsPerfect.exists_flat_model, Cohesive.isKFlat_twComp, Cohesive.exists_roof_sheafBC_restrOver, Cohesive.untwFRestrIso, Cohesive.isPerfect_untwF}
\leanok
Let D be resolving data on a compact space X, h ∈ 𝒜(X)\textsuperscript{2} closed and T a twist datum for h. Then every object K of D(TwMod(O, U, ȟ)) in D\textsubscript{perf}(X, ȟ) is isomorphic to J\textsubscript{T}(M) for some M in Ho(P\textsubscript{(𝒜(X), h)}). Steps proved on the way:

\begin{itemize}
\item every twisted perfect complex receives a quasi-isomorphism from a strictly bounded above complex K′ with componentwise flat terms, whose components K′\textsubscript{α} are then K-flat;
\item for such K′ (over pre-resolving data), F = untwist(K′) satisfies F|U\textsubscript{α} ≅ 𝒜\textsuperscript{0}|U\textsubscript{α} ⊗ K′\textsubscript{α} and is a perfect complex of 𝒜\textsuperscript{0}-modules on X;
\item base change along R → S of a bounded above flat complex preserves the existence of local roofs to strictly perfect complexes.
\end{itemize}
Deviation: K-flat variant (no flatness of 𝒜 over O assumed); X is not assumed Hausdorff here.
\par\smallskip{\footnotesize Source: blueprint-tw.md §4.}
\end{theorem}
\begin{proof}
\leanok
\uses{A11, A3, C7, S10, S3, S4, S8, S9, Z3.j, Z3.l, Z3.t}
\textbf{G.1} \texttt{D ≅ Q K}, bounded (C.6), \texttt{K' → K} componentwise flat (C.8).

\textbf{G.2} \texttt{F := glue(𝒜⁰ ⊗ K'\_α, u⁻¹(1 ⊗ φ))} ([A, Rem 7.2.4]).

\textbf{G.3} \texttt{F} globally bounded perfect on \texttt{X}.

\textbf{G.4} S10: \texttt{P̃ → F}, \texttt{P := Γ(P̃)}.

\textbf{G.6} A11 (general \texttt{c}, §1.2): \texttt{M}, \texttt{e ∈ Z0(M, N)}.

\textbf{G.8} \texttt{Ξ} glued, componentwise quasi-iso by C7.f.

\textbf{G.9} \texttt{ιT : K' → T\textasciicircum{}tw(K')} (C7.d).

\textbf{G.10} \texttt{Q(J\_T M) ≅ D}. \textbf{G′} (detail, 5 Nov 2026, for review at G5; refines G.1–G.10, no decision changed). Fix flat resolving data \texttt{P}, \texttt{h} closed of degree 2, a twist datum \texttt{T} (\texttt{h⁰\_α}, \texttt{u\_αβ}, \texttt{ȟ}), and \texttt{K} twisted perfect in \texttt{TwMod O U ȟ} with \texttt{D ≅ Q K}. \texttt{B\_α := 𝒜⁰|U\_α ⊗ −}, \texttt{η(k) = 1 ⊗ k} as in C.7′. \textbf{(a)} \emph{Flat model} (G.1): \texttt{K} has homology in \texttt{[a, b]} (C.6 (g)); \texttt{K' → K} a quasi-isomorphism from a bounded above complex with flat components (Z3.l, applied to \texttt{truncLE b}). \texttt{K'} is twisted perfect (C.6 (e)) and each \texttt{K'\_α} is K-flat on \texttt{Over U\_α} (S8.e). \textbf{(b)} \emph{Base change, untwisting} (G.2): \texttt{B := twBC ρ K'} (C.7′, cocycle \texttt{ρ(ȟ)}, components \texttt{B\_α(K'\_α) = F\_\{K'\_α\}} of C7.b by \texttt{rfl}). By B.5, \texttt{ρ(ȟ\_αβγ) = u\_βγ u\_αβ u\_αγ⁻¹}, so \texttt{g\_αβ := u\_αβ⁻¹} is a rescaling \texttt{ρ(ȟ) → 1} (\texttt{IsRescale}: \texttt{g\_αγ = g\_βγ g\_αβ ρ(ȟ\_αβγ)} ✓). \texttt{F := glue(rescale\_g B)}, a complex of \texttt{𝒜⁰}-modules on \texttt{X}, with gluing \texttt{φ\textasciicircum{}F\_αβ = u\_αβ⁻¹·φ\textasciicircum{}B\_αβ} and \texttt{φ\textasciicircum{}B(η k) = η(φ k)} (\texttt{bcφ\_bcUnit}). \texttt{F|U\_α ≅ F\_\{K'\_α\}} (C.7 (c)), compatible with the differentials. \textbf{(c)} \emph{Perfectness} (G.3): \texttt{F\_\{K'\_α\}} is perfect on \texttt{Over U\_α}: S8.g's \texttt{exists\_roof\_sheafBC} on over-sites, with C.7′ (a) in place of \texttt{sheafBCOverIso}. Hence \texttt{F} is perfect on \texttt{X} (local roofs through \texttt{F|U\_α ≅ F\_\{K'\_α\}} and \texttt{restr\_restrOver}), globally bounded (S9.d, \texttt{X} compact). \textbf{(d)} \emph{S10} (G.4): \texttt{Pt} strictly perfect over \texttt{𝒜⁰}, \texttt{f⁰ : Pt → F} a quasi-isomorphism; \texttt{P := Γ(Pt)} is a bounded complex of finitely generated projective \texttt{𝒜⁰(X)}-modules (S4.c) and \texttt{Γ(f⁰)} is a quasi-isomorphism (S3, \texttt{𝒜⁰} fine). \textbf{(e)} \emph{Local connections}: \texttt{θ\_α} the C7.b connection on \texttt{F\_\{K'\_α\}} over \texttt{Over U\_α} (\texttt{exists\_conn}; \texttt{θ\_α(η k) = 1 ⊗ η(d k)}), \texttt{T\_α := θ\_α.TC}, \texttt{ιT\_α : K'\_α → T\_α} a quasi-isomorphism (C7.d, \texttt{quasiIso\_ιT}, K-flat). \textbf{(f)} \emph{Twisted model} (G.9): \texttt{T\textasciicircum{}tw ∈ TwMod O U ȟ} (complexes) with components \texttt{T\_α} and gluing \texttt{a ⊗ s ↦ a ⊗ φ\textasciicircum{}B\_αβ(s)} on \texttt{Tot(𝒜(W), F\_\{K'\_α\}(W))} (\texttt{totMapAlong} with the identity of \texttt{𝒜(W)}; \texttt{φ\textasciicircum{}B} is \texttt{𝒜⁰(W)}-linear). Cocycle \texttt{ȟ} (acting through \texttt{ρ}). Chain maps: \texttt{θ\_β φ\textasciicircum{}B} and \texttt{φ\textasciicircum{}B θ\_α} are both \texttt{𝒜⁰}-connections, so their difference is \texttt{𝒜⁰}-linear out of \texttt{F\_\{K'\_α\}}, and it vanishes on \texttt{η k} (\texttt{d\_\{K'\}} commutes with \texttt{φ}), hence everywhere (C.7′ (b)). \texttt{ιT\textasciicircum{}tw : K' → T\textasciicircum{}tw} has components \texttt{ιT\_α} (commutes with the gluing: \texttt{φ\textasciicircum{}B η = η φ}) and is a quasi-isomorphism componentwise (C.6 (e)). \textbf{(g)} *Sections of the model of \texttt{F}* (C.7 (g)): for \texttt{V} open, the map \texttt{Tot(𝒜(V), F(V)) → ∏\_α Tot(𝒜(V ⊓ U\_α), F\_\{K'\_α\}(V ⊓ U\_α))}, restriction followed by the component \texttt{α} of \texttt{F|U\_α ≅ F\_\{K'\_α\}} (\texttt{π\_α}), is injective, and its image consists of the families with \texttt{y\_β = u\_αβ⁻¹·φ\textasciicircum{}T\_αβ(y\_α)} on the overlaps. Proof: \texttt{V ↦ Tot(𝒜(V), F(V))} is a sheaf (C7.a, for any \texttt{F}), and (b). \textbf{(h)} \emph{Curved connection} (G.5): for \texttt{s ∈ F\textasciicircum{}p(X)}, \texttt{θ\textasciicircum{}h(s)} is the element of (g) with components \texttt{θ\_α(s\_α) + h⁰\_α·s\_α}. Compatibility on \texttt{W ≤ U\_α ∩ U\_β}: on \texttt{s\_α = η k}, \texttt{θ\_β(u⁻¹η(φk)) + h⁰\_β u⁻¹η(φk) = u⁻¹(η(φ dk) + h⁰\_α η(φk))}, since \texttt{d(u\_αβ⁻¹) = −u\_αβ⁻¹(h⁰\_β − h⁰\_α)} (from \texttt{d u\_αβ = (h⁰\_β − h⁰\_α)u\_αβ}); in general both sides are connections (C.7′ (b)). \texttt{θ\textasciicircum{}h} is generator data over \texttt{(𝒜(X), h)} (A3.a): degree 1; flat, since in each component \texttt{(θ\_α + h⁰\_α)² = θ\_α² + d h⁰\_α = h|U\_α} (\texttt{h⁰\_α} odd); its form-degree-0 part is \texttt{d\_F} (\texttt{h⁰\_α} has form degree 1). \texttt{N\textasciicircum{}h := (F\textasciicircum{}•(X), θ\textasciicircum{}h)} is quasi-cohesive over \texttt{(𝒜(X), h)} with \texttt{(0,0)}-complex \texttt{Γ(F)}. Only \texttt{V = X} is needed: no curved copy of C7's \texttt{Conn} for all \texttt{V}. \textbf{(i)} \emph{Lifting} (G.6): A11 (general \texttt{c}) with \texttt{e⁰ := Γ(f⁰)} gives \texttt{M} cohesive over \texttt{(𝒜(X), h)} on \texttt{𝒜(X) ⊗ P} and \texttt{e ∈ Z0(M, N\textasciicircum{}h)} with \texttt{(0,0)}-component \texttt{e⁰} (curved copy of \texttt{exists\_cohesive\_gr0\_quasiIso}, steps (A)–(C)). \textbf{(j)} \emph{Localization} (G.7): \texttt{e\_α ∈ Z0(gauge\_\{h⁰\_α\}(M|U\_α), N\_α)}, \texttt{N\_α := θ\_α.N(U\_α)}, \texttt{x ↦ π\_α(extAlong(res) e (x))}. Closed: \texttt{π\_α ∘ res ∘ θ\textasciicircum{}h = (θ\_α + h⁰\_α) ∘ π\_α ∘ res} by (h), and the gauged differential is \texttt{𝕏 − h⁰\_α} (A.2). \textbf{(k)} \emph{Comparison} (G.8): \texttt{Ξ\textasciicircum{}tw : J\_T(M) → T\textasciicircum{}tw} with components \texttt{twLocalIso} (D.6, \texttt{W′ = U\_α}) followed by \texttt{Ξ\_\{e\_α\}} (C7.e). It commutes with the gluing: \texttt{J\_T}'s gluing is \texttt{u\_αβ·}, and \texttt{φ\textasciicircum{}T(π\_α y) = u\_αβ·π\_β(y)} by (g). Quasi-isomorphism componentwise by C7.f: \texttt{gr0Ξ\_\{e\_α\}} corresponds to \texttt{f⁰|U\_α} under \texttt{G0(gauge(M|U\_α)) ≅ Pt|U\_α} and \texttt{F|U\_α ≅ F\_\{K'\_α\}} (gauging changes only the form-degree-1 component), as in C5's step (D) with base \texttt{U\_α} in place of \texttt{⊤}. In detail: the generators of \texttt{gauge(M|U\_α)} are \texttt{𝒜⁰(U\_α) ⊗\_\{𝒜⁰(X)\} Γ(Pt\textasciicircum{}p)}, and \texttt{φ\_p : β ⊗ q ↦ β·q|U\_α} onto \texttt{Pt\textasciicircum{}p(U\_α)} is bijective because \texttt{Pt\textasciicircum{}p} is a summand of a finite free \texttt{𝒜⁰}-module (true for \texttt{𝒜⁰\textasciicircum{}m}, and preserved by retracts); \texttt{φ} intertwines \texttt{X0comp} with the differential of \texttt{Pt|U\_α} (gauging does not change \texttt{X0comp}). Then C5 (D) with base \texttt{U\_α} (\texttt{quasiIso\_gr0Ξ\_of\_equiv}, \texttt{Sc := Pt|U\_α}, \texttt{g := f⁰|U\_α} followed by \texttt{F|U\_α ≅ F\_\{K'\_α\}}) gives \texttt{gr0Ξ\_\{e\_α\}} a quasi-isomorphism; \texttt{hg} is \texttt{ε ∘ π\_α = π\_α ∘ ε} on \texttt{e(1 ⊗ q)} and the \texttt{(0,0)}-component \texttt{Γ(f⁰)} of \texttt{e} from (i). \textbf{(l)} \emph{Assembly} (G.10): \texttt{Q(J\_T M) ≅ Q(T\textasciicircum{}tw) ≅ Q(K') ≅ Q(K) ≅ D}. \textbf{(m)} \emph{Split and risks}: Z3.t = (b), (e)–(h), (j); Z3.u = (a), (c), (d), (i), (k), (l). Risks: (k) generalizing C5's step (D) (\texttt{quasiIso\_gr0Ξ\_abstract}, private, at \texttt{⊤}) to base \texttt{U\_α}; (g) the transport of \texttt{T(F)} sections along \texttt{F|U\_α ≅ F\_\{K'\_α\}}; curved/flat instance traps between \texttt{(𝒜(X), h)} and the flat \texttt{𝒜(U\_α)} in (h)–(j) (P15 notes; \texttt{cmpTot} as in \texttt{twCmp}).
\paragraph{Blueprint notes.}
G.1–G.4, G.6, G.8–G.10

Files: \texttt{Comparison/TwistedEssSurj*.lean}.
\end{proof}

\begin{theorem}[Z3.v — Main theorem twisted]
\label{Z3.v}
\lean{Cohesive.twistedBlockFunctor, Cohesive.MainStatementTwisted, Cohesive.main_theorem_twisted, Cohesive.essImage_twistedBlockFunctor}
\leanok
Let E be a finite-dimensional complex normed space and X a compact Hausdorff complex manifold modelled on E. Let h ∈ 𝒜\textsuperscript{0,2}(X) be ∂̄-closed, A = (𝒜\textsubscript{X}(X), ∂̄, h) the curved Dolbeault algebra and T a twist datum for h with cover (U\textsubscript{α}) and cocycle ȟ. Then the functor J\textsubscript{T} : Ho(P\textsubscript{A}) ⥤ D(TwMod(O\textsubscript{X}, U, ȟ)), M ↦ J\textsubscript{T}(M), into the derived category of ȟ-twisted sheaves of O\textsubscript{X}-modules is full and faithful, and its essential image is exactly D\textsubscript{perf}(X, ȟ) (the essential image statement holds for any choice of derived category).

Deviation: stated for each twist datum T (independence of T is not proved); D\textsubscript{perf} instead of D\textsuperscript{b\_}\{coh\}; twisted O\textsubscript{X}-modules only, curvature c = h for left modules, no factor 2πi.
\par\smallskip{\footnotesize Source: blueprint-tw.md §4.}
\end{theorem}
\begin{proof}
\leanok
\uses{Z3.s, Z3.u}
\textbf{H.1} \texttt{twistedComparison} for \texttt{CurvedResolvingData} (u).

\textbf{H.2} Dolbeault instance, \texttt{main\_theorem\_twisted}.

\textbf{H.3} K8: at \texttt{(0, trivialTwistDatum)} the source \texttt{HoP} is literally \texttt{HoP (dolbeault E X)} (P15); \texttt{glueEquiv : TwMod O \{⊤\} 1 ≌ Mod O}, \texttt{twistedBlockFunctor E X 0 T₀ ≅ blockFunctor ⋙ glue.mapDerivedCategory}-type isomorphism (\texttt{Equivalence.mapDerivedCategory} ≈ 40 lines from \texttt{mapDerivedCategoryCompIso}/\texttt{IdIso}/\texttt{NatTrans.mapDerivedCategory}), \texttt{TwPerf ↔ Perf}.
\paragraph{Blueprint notes.}
H.1–H.3

Files: \texttt{Comparison/TwistedEquivalence.lean}, \texttt{TwistedBlockFunctorDef.lean}, \texttt{StatementsTwisted.lean}, \texttt{MainTwisted.lean}, \texttt{ChecksTwisted.lean}.
\end{proof}
